% Part: turing-machines
% Chapter: machines-computations
% Section: introduction

\documentclass[../../../include/open-logic-section]{subfiles}

\begin{document}

\olfileid[fr]{tur}{mac}{int}
\olsection{Introduction}

Que signifie, pour une fonction de $\Nat$ dans
$\Nat$ par exemple, être \emph{calculable} ?
L'une des premières réponses, et la plus connue,
est qu'une fonction est calculable si une machine
de Turing peut la calculer. Alan Turing a proposé
cette notion en 1936. Les machines de Turing sont
un \emph{modèle de calcul} : elles donnent un sens
mathématiquement précis à l'idée de « procédure
de calcul ». La portée exacte de cette expression
fait débat, mais on admet largement qu'une machine
de Turing permet de spécifier une telle procédure.
Bien que l'expression « machine de Turing » évoque
un appareil physique pourvu de pièces mobiles,
une telle machine est, à strictement parler,
un objet purement mathématique. Elle idéalise
donc la notion de procédure de calcul. Nous
n'imposons, par exemple, aucune limite au temps
ni à la mémoire dont elle dispose : elle peut
calculer un résultat même si le calcul exigeait
plus d'emplacements de mémoire ou plus d'étapes
qu'il n'existe d'atomes dans l'univers.

\begin{explain}
Il est peut-être plus simple de concevoir une
machine de Turing comme le programme d'un mécanisme
imaginaire particulier. Ce mécanisme comprend un
\emph{ruban} et une \emph{tête de lecture-écriture}.
Dans notre version, le ruban est infini dans une
seule direction, vers la droite, et se divise en
\emph{cases}. Chacune peut contenir un symbole
appartenant à un \emph{alphabet} fini. Cet alphabet
peut comporter autant de symboles différents que
l'on veut, mais nous emploierons surtout les trois
symboles $\TMendtape$, $\TMblank$ et~$\TMstroke$.
Au démarrage, le ruban est vide, c'est-à-dire que
chaque case contient $\TMblank$, sauf la case la
plus à gauche, qui contient $\TMendtape$, et un
nombre fini de cases contenant l'\emph{entrée}.
À chaque instant, le mécanisme se trouve dans un
état pris parmi un nombre fini d'\emph{états}. Au départ, la tête
lit la première case à droite du marqueur de début
de ruban et la machine se trouve
dans un \emph{état initial} spécifié.\footnote{Correction éditoriale : l'introduction source place la tête sur la case du marqueur, mais la définition formelle et les exemples la placent dans la case suivante. Nous retenons cette seconde convention ; sinon la machine paire s'arrêterait immédiatement sur toute entrée.} À chaque étape,
le contenu de la case lue, l'état courant et le
programme déterminent ensemble l'action suivante.
Le programme est donné par une fonction partielle
qui, à un état~$q$ et un symbole~$\sigma$, associe
un triplet $\tuple{q', \sigma', D}$. Lorsque le
mécanisme est dans l'état~$q$ et lit~$\sigma$,
il remplace ce symbole par~$\sigma'$, déplace la
tête vers la gauche, vers la droite ou la laisse
sur place selon que $D$ vaut $\TMleft$, $\TMright$
ou $\TMstay$, puis passe dans l'état~$q'$.

Considérons par exemple la situation représentée
dans la \olref{fig:tm}.
\begin{figure}
  \olasset{\olpath/assets/diagrams/turing-machine.tikz}
  \caption{Une machine de Turing qui exécute son programme.}
  \ollabel{fig:tm}
\end{figure}
La partie visible du ruban porte le symbole de
début de ruban $\TMendtape$ dans la case la plus
à gauche, suivi de trois $1$, d'un $0$, puis de
quatre autres $1$. La tête lit la troisième case
à partir de la gauche, qui contient un~$1$, et
la machine est dans l'état~$q_1$ : nous disons
qu'elle « lit un $1$ dans l'état~$q_1$ ».
Si, pour l'entrée $\tuple{q_1,1}$, le programme
renvoie $\tuple{q_2,0,\TMstay}$, la machine
remplace ce~$1$ par un~$0$, laisse la tête sur
la même case et passe dans l'état~$q_2$.
Si le programme renvoie ensuite
$\tuple{q_3,0,\TMright}$ pour l'entrée
$\tuple{q_2,0}$, la machine réécrit le~$0$ par
un~$0$, sans changer en pratique le contenu de
la case, déplace la tête d'une case vers la
droite et passe dans l'état~$q_3$. Et ainsi
de suite.

La machine \emph{s'arrête} lorsqu'elle rencontre
un état~$q_n$ et un symbole~$\sigma$ pour lesquels
aucune instruction n'est prévue pour
$\tuple{q_n,\sigma}$, autrement dit
lorsque la fonction de transition n'est pas
définie sur $\tuple{q_n,\sigma}$. Elle n'a
alors plus d'instruction à exécuter. On peut
aussi représenter l'arrêt au moyen d'un état
spécial~$h$. Nous y reviendrons plus en détail.
\end{explain}

\begin{digress}
L'intérêt de l'article de Turing, « On computable
numbers », tient à ce qu'il présente une définition
formelle, mais aussi un argument selon lequel
cette définition saisit la notion intuitive de
calculabilité. Il paraît clair, d'après la
définition, qu'une fonction calculable par une
machine de Turing est calculable au sens intuitif.
Pour défendre la réciproque, à savoir que toute
fonction que nous jugerions naturellement calculable
l'est par une telle machine, Turing propose trois
types d'arguments, dans ses propres termes :
\begin{enumerate}
\item Un appel direct à l'intuition.
\item Une preuve de l'équivalence de deux définitions
  (au cas où la nouvelle définition aurait un attrait
  intuitif plus grand).
\item Des exemples de vastes classes de nombres
  calculables.
\end{enumerate}
Nous cherchons à définir la calculabilité « en
principe », sans tenir compte des limites pratiques
de temps et d'espace. Une fois cette notion générale
en place, on peut étudier les calculs qui disposent
de ressources bornées : c'est le cœur de la
\emph{complexité algorithmique}.
\end{digress}

\begin{history}
Alan Turing a introduit les machines de Turing en
1936. Il s'intéressait alors à la décidabilité de
la logique du premier ordre, mais son article est
également considéré comme fondateur pour la
conception des ordinateurs. Il y étudie surtout
les nombres réels calculables, dont on peut
produire le développement décimal, et remarque
que ses notions s'adaptent aux fonctions sur
les entiers naturels.

Cette publication précède de cinq ans le Z3
programmable de Konrad Zuse (1941), de sept ans
les premiers essais du Colossus conçu par Tommy
Flowers (fin 1943, avant sa mise en service à
Bletchley Park en 1944), de neuf ans l'achèvement
de l'ENIAC américain (1945), de douze ans le
prototype britannique Manchester Small-Scale
Experimental Machine (1948) et de treize ans les
premiers essais du BINAC américain (1949).
La machine expérimentale de Manchester fut la
première à exécuter un programme conservé dans
une mémoire électronique. À la différence de
nombre de machines antérieures, elle conservait
ainsi ses instructions en mémoire, sans qu'il
faille reconfigurer physiquement son câblage
pour chaque nouveau programme.\footnote{Corrections historiques : nous distinguons le Z3 de 1941 des machines antérieures de Zuse, attribuons Colossus à Tommy Flowers et distinguons ses essais de fin 1943 de son fonctionnement à Bletchley Park en 1944. La priorité de Manchester porte sur un programme conservé en mémoire électronique ; les machines précédentes n'exigeaient pas toutes un recâblage. Voir le National Museum of Computing, \url{https://www.tnmoc.org/colossus}, et le Computer History Museum, \url{https://www.computerhistory.org/timeline/1941/} et \url{https://www.computerhistory.org/revolution/story/87}.}
\end{history}

\end{document}
